\documentclass[11pt]{article}

%% Packages
\usepackage[T1]{fontenc}
\usepackage[utf8]{inputenc}
\usepackage[left=0.5in, right=0.5in, top=0.35in, bottom=0.35in]{geometry}
\usepackage{enumitem}
\usepackage[hidelinks]{hyperref}
\usepackage{titlesec}
\usepackage{amssymb}

%% Page style
\pagestyle{empty}
\setlength{\parindent}{0pt}
\setlength{\parskip}{0pt}

%% Section formatting
\titleformat{\section}
  {\bfseries\uppercase}
  {}{0em}{}
  [\vspace{1pt}\hrule\vspace{2pt}]
\titlespacing*{\section}{0pt}{4pt}{1pt}

%% Compact bullets
\setlist[itemize]{
  noitemsep, topsep=0pt, parsep=0pt,
  partopsep=0pt, leftmargin=1.2em, label=\textbullet
}

%% Entry header — optional [url], {Title/Org}, {Role/Tech}, {Date}
\newcommand{\entryhead}[4][]{%
  \noindent
  \ifx&#1&%
    \textbf{#2}%
  \else
    \href{#1}{\textbf{#2}}%
  \fi
  {\small\textit{ $\cdot$ #3}}%
  \hfill {\small #4}%
}

% ─────────────────────────────────────────────
\begin{document}

%% ── HEADER ──────────────────────────────────
\begin{center}
  {\LARGE\textbf{Karthik Shivashankaran}} \\[2pt]
  {\small
    Brooklyn, NY \textbar{}
    +1 848-702-5284  \textbar{}
    \href{mailto:ks7926@nyu.edu}{ks7926@nyu.edu} \\
    \href{https://imaginelenses.com/}{imaginelenses.com} \textbar{}
    \href{https://www.linkedin.com/in/imaginelenses}{linkedin.com/in/imaginelenses} \textbar{}
    \href{https://github.com/imaginelenses}{github.com/imaginelenses}
  } \\[1pt]
  {\small\textbf{US Work Authorized --- OPT (STEM Extension Eligible)}}
\end{center}

\vspace{-4pt}

%% ── SUMMARY ─────────────────────────────────
\section*{Summary}
M.S.\ Computer Engineering (2026) specializing in Robotics and Autonomous Navigation. Built and deployed full autonomy stacks on physical hardware, from perception through planning; implemented state estimation and control algorithms from scratch. Open-source creator with 130+ GitHub stars.

%% ── EDUCATION ───────────────────────────────
\section*{Education}
\textbf{New York University} --- M.S.\ Computer Engineering $\cdot$ Scholarship Recipient
\hfill New York, NY \textbar{} 2026 \\
{\small\textit{Courses: Robot Localization \& Navigation, Reinforcement Learning for Robotics, Computer Vision, , Real-Time Embedded Systems, Computer Architecture, Operating Systems}}

\smallskip
\textbf{Acharya Institute of Technology} --- B.E.\ Mechatronics
\hfill Bangalore, India \textbar{} 2022

%% ── TECHNICAL SKILLS ────────────────────────
\section*{Technical Skills}
{\small
\textbf{Languages:} Python, C, C++, Java, JavaScript, SQL, CMake, MATLAB \\[1pt]
\textbf{ML / Computer Vision:} YOLOv8, Object Detection, Multi-Object Tracking, ByteTrack,
  3D Gaussian Splatting, Depth \& Pose Estimation, Structure-from-Motion (SfM),
  Camera Calibration, PyTorch, OpenCV, NumPy \\[1pt]
\textbf{Robotics \& Autonomy:} SLAM, Motion Planning, Path Planning, Sensor Fusion,
  State Estimation (Kalman Filter, EKF, Particle Filter), Open3D, Point Cloud Processing,
  Reinforcement Learning (Stable-Baselines3), Gymnasium, ROS2 (familiar) \\[1pt]
\textbf{Control Systems:} PID, MPC, LQR, MPPI, Closed-Loop Control, Real-Time Systems \\[1pt]
\textbf{Embedded \& Hardware:} ESP32, Arduino Giga R1, Raspberry Pi, IMU, Circuit Design,
  CAD, 3D Printing \\[1pt]
\textbf{Systems \& Tools:} Git, Docker, CI/CD, Linux,
  AWS (EC2, S3, Lambda), TCP/IP, WebSockets, CAN, UART, SPI, I2C
}

%% ── APPLIED RESEARCH & TECHNICAL PROJECTS ────────────────────────────────
\section*{Applied Research \& Technical Projects}
\entryhead[https://github.com/imaginelenses/pixelsToSteps]
{pixelsToSteps}
{Python $\cdot$ NumPy $\cdot$ SciPy $\cdot$ OpenCV $\cdot$ Gymnasium}
{May 2026}
\begin{itemize}
  \item Achieved 100\% survival across 100 trials under true plant mismatch (pole mass $\times$1.1, half-length $\times$0.9) stabilizing a cartpole from raw pixels alone; used an LQR teacher and hybrid Luenberger observer (Ridge regression, 40 demos), up from 96/100 at half the data
  \item Observer maps Gaussian-blur + Otsu-thresholded frames to pole angle (R$^2$\,=\,0.987, RMSE\,=\,0.123\textdegree{}), fusing pixel estimates with linearized dynamics at 60\,Hz as a single matrix multiply with no angle sensor at inference
\end{itemize}

\vspace{1pt}
\entryhead[https://github.com/imaginelenses/NYU_ROB_GY_6213]
  {Autonomous Robot Navigation Stack}
  {Python $\cdot$ C++ $\cdot$ UDP $\cdot$ RPLidar $\cdot$ ArUco $\cdot$ SciPy}
  {Apr 2026}
\begin{itemize}
  \item Closed the full autonomy loop on a physical Ackermann robot (Arduino Giga R1 + ESP32 for motor control and sensing, compute offloaded to host over UDP) with a from-scratch MPPI controller; sampled thousands of stochastic trajectories per step, weighted by an EDT-based composite cost (goal proximity + obstacle clearance), achieving real-time autonomous navigation
  \item Fused UDP-streamed wheel encoder odometry (10\,Hz) with ArUco visual pose measurements via a from-scratch EKF, reducing positional RMSE to 4.14\,cm and recovering from pose errors in $<$0.1\,s
  \item Built a 3,000-particle Monte Carlo Localizer for global localization against a pre-built LiDAR occupancy map; achieved 0.27\,cm position spread and 0.90\textdegree{} heading $\sigma$, recovering from a 1\,m\,/\,90\textdegree{} initialization error in 0.63\,s
\end{itemize}

\vspace{1pt}
\entryhead[https://github.com/exploring-curiosity/2Pass-3D-Reconstruction-of-Urban-Scene]
  {3D Line-of-Sight Audit System for Intersections}
  {Python $\cdot$ Gaussian Splatting $\cdot$ OpenCV}
  {Dec 2025}
\begin{itemize}
  \item Owned the 3D reconstruction pipeline within a 5-person team-built intersection audit system, replacing manual checks with an 8-camera SfM and Gaussian Splatting pipeline (30 FPS); generated a 1.34M-point cloud to evaluate pedestrian visibility
  \item Contributed to deployment of YOLOv8 + ByteTrack for real-time multi-object tracking (20--40 tracks/ camera); multi-camera fusion achieved ${>}$90\% precision, ${>}$85\% recall, and detected vehicles ${\sim}$12.2s earlier than a single-camera 2D baseline
\end{itemize}

\vspace{1pt}
\entryhead[https://smartswitch.imaginelenses.com/]
  {Smart Switch}
  {JavaScript $\cdot$ YAML $\cdot$ ESP32 $\cdot$ Circuit Design $\cdot$ CAD $\cdot$ 3D Printing}
  {Aug 2023}
\begin{itemize}
  \item Owned end-to-end from circuit design and 3D-printing through firmware and auto-generated YAML config tooling; eliminated cloud round-trip latency (sub-ms LAN vs.\ 80--200\,ms cloud), 5 units deployed and operational
\end{itemize}

\vspace{1pt}
\entryhead[https://blendit.imaginelenses.com/]
  {Blendit \normalfont{\small($\bigstar$ 130 GitHub Stars)}}
  {Python $\cdot$ Git $\cdot$ Blender}
  {Sep 2022}
\begin{itemize}
  \item Git-based version control extension for Blender; reduced project file size by 83\% via function-call delta logging instead of full binary snapshots
\end{itemize}

\end{document}    